\documentclass[11pt]{article}
\usepackage[utf8]{inputenc}
\usepackage[letterpaper,top=1in,right=1in,left=1in,bottom=1in]{geometry}
\usepackage{amsmath}
\usepackage{amsfonts}
\usepackage{amssymb}
\usepackage{graphics}
\usepackage[demo]{graphicx}
\usepackage{placeins}
\usepackage[capposition=top]{floatrow}
\usepackage{subcaption}
\usepackage{lscape}
\usepackage{pdflscape}
\usepackage{setspace}
\usepackage{standalone}
\usepackage{fancyhdr}
\usepackage{framed}
\usepackage{color}
\usepackage{tikz}
\usepackage{pgfplots}
\usepackage{longtable}
\usepackage{booktabs,caption}
\setlength{\marginparwidth}{2cm}
\usepackage[draft]{todonotes}
\usepackage{booktabs,caption}
\usepackage[flushleft]{threeparttable}
\usepackage{bbm}
\usepackage{multirow}
\usepackage{lineno}
\usepackage{titlesec}
\usepackage{hyperref}
\usepackage{comment}
\usepackage[draft]{todonotes}
\usepackage[justification=centering]{caption}
\pgfplotsset{compat=1.15}
\newtheorem{theorem}{Theorem}[section]
\newtheorem{assumption}{Assumption}[section]
\usepackage{longtable}
\usepackage{array}

\usepackage[authoryear]{natbib}
\bibliographystyle{ecta}

\usepackage{colortbl}
\usepackage{ragged2e}
\usepackage{xfp}

\usepackage{stackengine}
\usepackage{array}
\newcolumntype{L}[1]{>{\raggedright\arraybackslash}p{#1}}
\setstackEOL{\#}
\setstackgap{L}{12pt}

\usepackage[page,toc,titletoc,title]{appendix}
\usepackage{blindtext}
\newcommand{\unitlabel}{standard deviations}
\usepackage{pdfpages}

\hypersetup{
	colorlinks=true, %set true if you want colored links
	linktoc=all,     %set to all if you want both sections and subsections linked
	linkcolor=black,  %choose some color if you want links to stand out
	citecolor=blue
}

%\linenumbers 

\title{Cross-grade ability grouping and education inequality \\ Evidence from Kenya and Liberia}

\begin{document}
\author{Mridul Joshi}
\date{Please do not circulate \\ \today}

\maketitle
\doublespacing

\section{Introduction}

Should students be grouped to instruction based on the number of years of education they have received or based on their mastery of material? In the standard model, level of instruction is determined, in large part, by the quantity of education that a student has received. However, other characteristics of a student may allow for more efficient targeting of instruction. If tests can more accurately forecast learning needs than years of education, then grouping students across grades on the basis of tests may result in more efficient targeting of instruction. Recent evidence that aligning instruction to students' level can result in more rapid learning \citep{banerjee/cole/duflo/linden:2007,duflodupaskremer,banerjee2017proof,duflo2020experimental,muralidharan2019disrupting} suggests that if the distribution of learning levels in adjacent grades overlaps substantially, schools effectiveness could be improved by grouping children on the basis of past performance. 

Tracking on the basis of a prior test score raises several concerns. First, tracking students raises distributional concerns. Students who are assigned to lower levels of instruction may be harmed either through adverse peer effects or by reduced self-esteem \citep{betts:2011}. Second, there is a risk of mis-classification. Test scores are noisy measures of ability, and even a high quality test may not capture the domain that is most relevant to targeting instruction. A test that is a noisy measure of the skills needed to benefit from a given level of instruction may ultimately lead to less targeted instruction. 

This study examines the impact of a reading group intervention in schools operated by NewGlobe Education in Kenya and Liberia. In Kenya, the program tracked students in grades 1 and 2 into one of two reading groups based on performance on literacy and English tests. In Liberia, the program tracked students in Grades 1 through 4, where students in Grade 1 and 2 were potentially grouped together, as were students in Grades 3 and 4.

The notion of grouping children by past academic performance instead of age or quantity of education has a long history. Grouping children by level of achievement instead of age is an increasingly popular idea both rich and developing countries. Programs grouping children between Kindergarten and third grade are increasingly being experimented with in schools in the United States \citep{edweek20191106}. Cross-grade ability grouping for literacy instruction was proposed as a way to improve the efficiency of education in Joplin, Missouri in the mid-20th century \citep{floyd:1954}. The Joplin Plan, as it came to be known, was adopted by several school districts in the 1950s and 60s. \cite{slavin:1987} reviewed evidence on tracking programs, and found that cross-grade ability grouping for reading only was among the most effective tracking interventions. 

[INCOMPLETE]The present study contributes to this literature by providing the first large-scale experimental evidence from low- and middle-income countries (LMICs). We provide evidence from two separate experiments, one in public schools in Liberia, one of the poorest countries in the world and a setting where learning levels are very low, and the other from private schools in Kenya, a lower-middle income country with one of the most effective education systems in Sub-Saharan Africa. The intervention was specifically motivated by studies showing the benefit of differentiated instruction in schools in low- and middle-income countries. The unique setting allows for greater control over the instructional level than previous studies. This feature allows for direct evaluation of the effects of mismatched instructional level. 

RESULTS PARAGRAPH

SHORT DISCUSSION PARAGRAPH

ROADMAP PARAGRAPH

%Tracking interventions typically rely on student test scores to determine which ``reading group" students should be assigned to. While reading group interventions may be effective overall, there is a risk of misclassification, and in those cases, students may receive inappropriate instructional levels. This study will explore the risk of misclassification by identifying students for whom the score used for tracking my contain important measurement error. First, test scores of students who enter Bridge improve rapidly over the first year of their exposure to the new schools. While there are many potential explanations for this, what is important for the problem of misclassification is that, for students who have entered Bridge recently and are improving rapidly, the test score is less informative about what the student's ability likely is at endline. To explore whether this type of misclassification affects the intervention, I compare the performance of students who have entered recently to students who have been in Bridge for several years. These results indicate that the 

\section{Literature review}
The question of whether pupils in a given grade benefit from being tracked into groups based on their initial achievement\footnote{There is also a large body of research on between-school ability grouping, mainly from Europe where tracking between academic and vocational tracks is prevalent. In this paper, we are interested in within-school ability grouping.} is a recurring question in the economics of education literature. Proponents of tracking posit that since it reduces the variation in achievement among pupils, tracking allows teachers to target instruction to the level of the pupils, thus benefiting students in both the higher and the lower tracks. On the other hand, critics of tracking argue that tracking disproportionately hurts the already low-achieving students through negative peer effects as well as loss of self-esteem through knowledge of being placed in the lower track. 

Despite its long history, the evidence around the effectiveness of ability grouping is rather mixed (see \cite{betts:2011}) for a comprehensive review). Most of the evidence on within-school tracking comes from high income countries, primarily the United States and hints towards a zero impact of tracking. \citep{slavin:1987} does not find a statistically significant impact of tracking in his meta-analysis of 14 studies on tracking in elementary schools, of which one had random assignment (cite xx) and three were large-scale, longitudinal studies (cite xx). At the secondary school level, none of the six randomised studies reviewed in \cite{slavin1990achievement} found a positive effect of tracking, although there was some evidence of a small negative effect \citep{hanushek2006does}. Other recent studies and reviews find that ability grouping has little to no impact on student achievement \citep{kulikand1987effects, terrin2022effect}.

Given the nature of tracking, there is likely to be heterogeneity in effects between high and low ability tracks. The overall evidence is inconclusive but usually suggests that tracking may be beneficial for students assigned to the high track and harmful for students assigned to the low track \citep{fu2018ability, lavrijsen2015new, burke2013classroom, hoxby2005taking}. Other studies report large, positive effects of tracking on academic achievement for students in the higher track and no negative impact on lower track students \citep{card2016can, vardardottir2013peer, antonovics/black/cullen/meiselman:2022}. On the other hand, \cite{a} Antecol et al (2013) found that tracking benefits the low ability track students and hurts the high achievement class. On the other hand, \cite{garlick} finds that tracking (in university dormitories) reduces the GPA of low-scoring students but does not impact high-scoring students. Albeit in the context of higher education, a recent study by \cite{de2022college} employs a regression discontinuity design and reports that students who were marginally admitted to the higher ability class in a public university in Columbia were less likely to pass. 

Well-identified studies on tracking or ability grouping are virtually absent in developing countries.  Among the few, \cite{duflodupaskremer} find positive effects of tracking in schools in Kenya. Students in tracking schools scored 0.18 standard deviations higher than students in non-tracking schools that persisted after the programme ended. The effects were observed on all quantiles of the baseline ability distribution.\footnote{\cite{cummins2017heterogeneous} reanalyzes the data from  \cite{duflodupaskremer} show that there is a positive effect in the low track only when the class is assigned to a contract teacher.} Apart from this, most other studies comprise ‘partial tracking’ in the spirit of Teaching at the Right Level (TaRL) and find significant gains to partial tracking (Banerjee et al, 2017; Duflo et al, 2022). 

Tracking interventions effectively change three components - the level of instruction, the dispersion in classroom achievement as well as the mean classroom achievement.  \cite{duflodupaskremer}, postulate that tracking allows teachers to teach at the level of the students and therefore positively impact both high ability and low ability students. They show that low ability students gain basic knowledge skills while high achieving students gain advanced skills suggesting that tracking allowed teachers to tailor their instruction to the level of the students. Another study by Evertson, Sanford, and Emmer (1981) supports this conclusion.  

The evidence on the impact of dispersion in the classroom as well the mean classroom achievement is rather mixed. There is evidence that suggests that being surrounded by more able peers leads to an increase in achievement (add cites). However, there are many empirical studies that indicate that classroom peer effects may be neglible (was it Angrist? cites check) or even negative ( Prof. Hoxby's invidious comparisons - see Luc Beheghel's notes).

Studies on classroom inputs often overlook the effects of dispersion in classroom achievement in favour of average classroom achievement (Rjosk, 2022). Kuzmina and Ivanova (2018) show that dispersion in achievement is positively correlated with reading outcomes. Similarly, Vigdor and Nechyba (2007) found a positive relation between dispersion in classroom achievement and student outcomes. Other studies that document that positive relationship are Scharenberg (2012) and  Chiu et al. (2017). On the other hand, there are handful of studies that find a negative association  (Duflo et al, 2011; Ding and Lehrer, 2007) or no association (Hanushek et al, 2003) between dispersion in peer achievement and student outcomes

Add paragraph on educational inequality due to tracking - tons of work on this - I have a list of cites somewhere.  

Our study also adds to the literature of classroom mixing. The literature provides no consensus on the effect of mixed grade classrooms. A survey of 56 studies concluded that there are no consistent differences in student outcomes between multi-grade and single grade classrooms.  Sims (2008) shows that mixed-grade classrooms are associated with a reduction in test scores for pupils in grade 2 and grade 3 assuming negligible class size effects. On the other hand, Leuven and Ronning (2011) show that students in mixed grades outperform pupils in single grade classroom in both high stakes exams as well as in tests set by their teachers. The authors conclude that whether the overall effect is positive or negative or zero, depends on the type of exposure to high and low ability peers. 

Our intervention at hand is a type of `cross-grade ability grouping' or `regrouping' \citep{slavin:1987} that has been termed as the Joplin Plan \citep{floyd:1954} after the city in Missouri where it was first piloted. The Joplin Plan creates mixed classrooms by regrouping students for reading based on their reading achievement instead of their grades or age. Once the reading class is over, students return to their original grades for other subjects. Multiple studies have tried to evaluate the impact of the Joplin Plan. Early studies found mixed results, although the research designs often did not include a sufficient number of randomized units be convincing one way or another. Most studies in the meta-analysis by \cite{slavin:1987} report a positive impact (Morgan and Strucker, 1960; Hilsen et al, 1967), there are some that report a negative impact (Kierstead, 1963) or a zero impact \citep{carson/thompson:1964,powell:1964}. However, according to Hiebert (1987) the overall positive effect of the Joplin Plan documented in \cite{slavin:1987} rests on assumptions about the nature of the interventions studied and may not necessarily hold up. 

To the best of our knowledge, ours is the first study that uses random assignment to test the impact of the Joplin Plan in two developing country settings. 

Misclassification has been noted as a potential problem for tracking \citep{betts:2011} [need more cites for this]. Add paragraph using Piopiunik (2013), Brunello, Giannini and Ariga (2007), Bedard (1997), Allen and Barnsley (1993) - maybe Bergman et al (2021) has something regarding prediction analytics for tracking. Also check Diris (2012)

Tracking changes the peer groups that students are exposed to. One of the primary concerns raised against tracking is that students tracked into lower achieving groups may be harmed by exposure to lower performing peers \citep{betts:2011}. 


\section{Background \& Design}

This study takes place in two separate countries, in two very different school systems.

Kenyan education is a regional exemplar in the education space, although learning levels remain below those of wealthy countries. 

Liberia is one of the poorest countries on earth. Although in recent years it has grown rapidly and improved access in education, Liberia remains a regional laggard. 

Class sizes in Kenya are pretty small relative to those in public schools in Kenya and in the region.

Class sizes in Liberian public schools operated by Bridge are similar to those seen elsewhere in sub-Saharan Africa.

\subsection{Experiment 1 (Kenya)}

In August 2018, Bridge selected 50 academies out of the 300 academies that it was operating in Kenya to receive the reading group intervention.  The 300 academies represented nearly all of the schools currently being operated in Kenya at that time. 

In June 2018, prior to assignment, baseline literacy and English tests were administered in 292 schools, including all 50 schools in the treatment. In treated academies, students in Standards 1 and 2 were assigned to either a Blue or Purple reading group. In treated schools, Standard 1 students were assigned to the Purple group if they scored above 40 total points on a combined test of English and Literacy, whereas Standard 2 students were assigned to the Purple group if they scored above a 35. %The test used for assignment was different in each grade and the tests administered at endline are different, so it is not possible to say whether this choice actually led to lower ``ability'' students being assigned to the purple group in STD 2.

After assignment, students in treated schools then received 90 minutes of instruction with their color group each day for Term 3 (approximately 3 months) of 2018. In control schools, students received the status quo: 90 minutes of instruction according to grade.

Treatment assignment was stratified at the province level. 

The intervention lasted from September until the end of the 2018 school year in December. 

Timeline of events in both studies.
Consort diagram for both studies.

\subsection{Identification strategy}

Say we are interested in the effect of average peer quality on student achievement. Usual reasons why we cannot run an OLS regression of student outcome on mean peer outcome. We can leverage a random variation in the peers of students due to a tracking treatment. However, tracking is conditionally random because it also depends on your pretest ability: schools are first assigned to a treatment or control condition and students in the treatment group are assigned to high or low tracks based on the pretest ability. 

Following Hoxby and Weingarth (2005), we can think of each student's simulated cohort i.e. the peers they will have based on the assignment mechanism. We can compute means of pretest scores using this simulated cohort as an instrument for peer quality. 

This will potentially lead to high ability students getting better high ability peers and low ability students getting low ability peers deterministically. One way to deal with this issue is to condition on the expected value of treatment to eliminate this selection bias by effectively controlling for the non-random variation in treatment (Imbens and Hirano, 2004). This expected value is easy to compute because we know the assignment mechanism. 

\singlespacing
\bibliography{bib}

\end{document}